\subsection{Main Results at Standard Annotation Budget}
\label{sec:main_results}

\cref{tab:main} presents the primary comparison at $n_\text{rare} = 20$ labeled rare-class training cells. Results are mean $\pm$ SD over three random seeds (single-seed entries indicate that only one complete run was available at time of writing).

\begin{table}[t]
\centering
\small
\caption{Rare-class F1 across datasets and methods ($n_\text{rare}=20$, mean $\pm$ SD over 3 seeds unless noted). \textbf{Bold} marks the best result per row. Dashes (--) indicate baselines not yet evaluated for this dataset. LR = logistic regression (CellTypist-equivalent).}
\label{tab:main}
\begin{tabular}{llccccc}
\toprule
Dataset & Rare class & scANVI & kNN & LR & \textbf{Gate+Marker} & Fusion-gated \\
\midrule
Immune DC & ASDC & $0.656{\pm}0.147$ & $0.614{\pm}0.015$ & $0.761$ & \textbf{$0.933{\pm}0.019$} & $0.931{\pm}0.023$ \\
 & cDC1 & $0.208{\pm}0.239$ & $0.807{\pm}0.086$ & $0.977$ & \textbf{$0.985{\pm}0.002$} & $0.982{\pm}0.003$ \\
\midrule
Human pancreas & $\varepsilon$-cell & \textbf{$0.889{\pm}0.192$} & $0.444{\pm}0.385$ & \textbf{$1.000$} & $0.889{\pm}0.192$ & $0.722{\pm}0.255$ \\
 & $\gamma$-cell & \textbf{$0.996{\pm}0.003$} & $0.805{\pm}0.129$ & $0.903$ & $0.992{\pm}0.009$ & $0.994{\pm}0.006$ \\
\midrule
Tabula Liver & NCM & $0.374{\pm}0.116$ & $0.416{\pm}0.179$ & $0.313$ & $0.625{\pm}0.115$ & \textbf{$0.704{\pm}0.119$} \\
Tabula Pancreas & $\beta$-cell & $0.897{\pm}0.070$ & \textbf{$0.913{\pm}0.067$} & $0.966$ & $0.897{\pm}0.070$ & $0.897{\pm}0.070$ \\
Tabula Spleen & ILC & $0.625$ & -- & $0.653$ & \textbf{$0.857$} & $0.764$ \\
Tabula Kidney & Endothelial & \textbf{$0.928{\pm}0.038$} & -- & -- & $0.919{\pm}0.057$ & \textbf{$0.928{\pm}0.038$} \\
PBMC & pDC & \textbf{$0.999$} & -- & -- & \textbf{$0.999$} & \textbf{$0.999$} \\
\bottomrule
\end{tabular}
\end{table}

The results confirm a clear pattern organised by separability regime. For the four rare cell types with $\sepratio > 1.3$ \emph{and} scANVI baseline F1 below 0.75 (ASDC, cDC1, NCM, ILC), Gate+Marker achieves average F1 gains of $+38.4$ percentage points over scANVI (range $+23.2$ to $+77.7$ pp). For the two cell types with high separability but already-reliable scANVI (endothelial cell, $\sepratio = 1.69$, baseline F1 $= 0.928$; pDC, $\sepratio = 2.03$, baseline F1 $= 0.999$), Gate+Marker produces near-zero modification (mean $|\Delta\text{F1}| < 0.01$), correctly recognizing that there is nothing to rescue. For the three cell types with $\sepratio < 1.2$ ($\varepsilon$-cell, $\gamma$-cell, $\beta$-cell), the module falls back to the scANVI prediction with a mean absolute F1 change of $0.003 \pm 0.007$.

Notably, logistic regression (LR) outperforms scRareRefine on $\varepsilon$-cells (LR F1 $= 1.000$ vs. Gate+Marker $= 0.889$) and $\beta$-cells (LR $= 0.966$ vs. $0.897$), both low-separability cases where scANVI is unreliable and expression-space features provide additional discriminability. In contrast, on the four high-separability cases where scRareRefine is designed to operate, it outperforms LR by an average of $+15.9$ pp (ASDC: $+17.2$ pp, cDC1: $+0.8$ pp, NCM: $+31.2$ pp, ILC: $+20.4$ pp), confirming that the separability ratio predicts the regime in which prototype-based rescue dominates expression-space classifiers.

\subsection{Data Efficiency under Extreme Annotation Budgets}
\label{sec:efficiency}

\cref{fig:efficiency} and \cref{tab:asdc_ablation,tab:cdc1_ablation} show how F1 varies with rare training size $n_\text{rare} \in \{5, 10, 20, 50, \text{all}\}$.

\begin{figure}[t]
\centering
\includegraphics[width=\linewidth]{fig_data_efficiency_v2.pdf}
\caption{Rare-class F1 as a function of training size for ASDC (left) and cDC1 (right) in the Immune DC dataset. Lines are means over three seeds; shaded bands are $\pm 1$ SD. Gate+Marker and Fusion-gated maintain high F1 even at $n_\text{rare} = 5$, where scANVI and kNN nearly or completely fail.}
\label{fig:efficiency}
\end{figure}

For cDC1, the headline result is at $n_\text{rare} = 5$: scANVI achieves F1 $= 0.003 \pm 0.005$, kNN achieves F1 $= 0.000 \pm 0.000$, and Gate+Marker achieves F1 $= 0.986 \pm 0.004$. The improvement of $+0.983$ F1 points at this annotation budget is consistent across all three seeds, demonstrating that the rescue pipeline extracts reliable rare-cell signal from latent geometry alone when only five labeled examples are available in training.

For ASDC, Gate+Marker achieves F1 $= 0.904 \pm 0.011$ at $n_\text{rare} = 5$ versus $0.025 \pm 0.031$ for scANVI and $0.000$ for kNN ($+0.879$ gain). As $n_\text{rare}$ increases, the scANVI baseline catches up: at $n_\text{rare} = 50$, scANVI reaches $0.879$ and Gate+Marker reaches $0.924$, a reduced but still positive margin.

For the pancreatic $\varepsilon$-cell (moderate separability, $\sepratio \approx 1.07$), Gate+Marker provides an F1 gain of $+0.22$ at $n_\text{rare} = 5$ (from scANVI F1 $\approx 0.00$ to Gate+Marker $\approx 0.22$), but converges to baseline performance by $n_\text{rare} = 20$. This illustrates Regime~2 from \cref{sec:problem}: even at moderate separability, the pipeline provides additional value when the annotation budget is so low that the scANVI classifier has insufficient data to form a reliable decision boundary.

\begin{table}[h]
\centering
\small
\caption{ASDC rare-class F1 vs.\ training size (mean $\pm$ SD over 3 seeds, immune DC, batch-heldout split)}
\begin{tabular}{lrrrr}
\hline
$n_\text{rare}$ & scANVI Baseline & kNN ($k$=15) & \textbf{Gate+Marker} & \textbf{Fusion-gated} \\
\hline
5 & $0.025 \pm 0.031$ & $0.000$ & \textbf{$0.904 \pm 0.011$} & $0.899 \pm 0.017$ \\
10 & $0.162 \pm 0.227$ & $0.250 \pm 0.110$ & $0.870 \pm 0.091$ & \textbf{$0.913 \pm 0.016$} \\
20 & $0.656 \pm 0.147$ & $0.614 \pm 0.015$ & \textbf{$0.933 \pm 0.019$} & $0.931 \pm 0.023$ \\
50 & $0.879 \pm 0.043$ & $0.805 \pm 0.026$ & $0.924 \pm 0.002$ & \textbf{$0.930$} \\
All & $0.947 \pm 0.003$ & $0.924 \pm 0.007$ & \textbf{$0.951 \pm 0.007$} & $0.951 \pm 0.003$ \\
\hline
\end{tabular}
\label{tab:asdc_ablation}
\end{table}

\begin{table}[h]
\centering
\small
\caption{cDC1 rare-class F1 vs.\ training size (mean $\pm$ SD over 3 seeds, immune DC, batch-heldout split)}
\begin{tabular}{lrrrr}
\hline
$n_\text{rare}$ & scANVI Baseline & kNN ($k$=15) & \textbf{Gate+Marker} & \textbf{Fusion-gated} \\
\hline
5 & $0.003 \pm 0.005$ & $0.000$ & \textbf{$0.986 \pm 0.004$} & $0.985 \pm 0.003$ \\
10 & $0.034 \pm 0.030$ & $0.292 \pm 0.221$ & \textbf{$0.979 \pm 0.012$} & $0.977 \pm 0.011$ \\
20 & $0.208 \pm 0.239$ & $0.807 \pm 0.086$ & \textbf{$0.985 \pm 0.002$} & $0.982 \pm 0.003$ \\
50 & $0.992 \pm 0.004$ & $0.968 \pm 0.007$ & \textbf{$0.997 \pm 0.001$} & $0.992 \pm 0.004$ \\
All & $0.992$ & \textbf{$0.996 \pm 0.002$} & $0.991 \pm 0.002$ & $0.989 \pm 0.005$ \\
\hline
\multicolumn{5}{l}{$\dagger$ fewer than 3 seeds available} \\
\end{tabular}
\label{tab:cdc1_ablation}
\end{table}


\subsection{Separability Ratio as a Pre-deployment Diagnostic}
\label{sec:sepratio_result}

\cref{fig:sepratio} shows rare-class F1 gain (Gate+Marker minus scANVI) plotted against the cell-type-level mean separability ratio at $n_\text{rare} = 20$, aggregating over three seeds.

\begin{figure}[t]
\centering
\includegraphics[width=0.65\linewidth]{fig_separability_gain.pdf}
\caption{Separability ratio versus F1 gain (Gate+Marker $-$ scANVI) at $n_\text{rare} = 20$. Each point represents a rare cell type; error bars show $\pm 1$ SD across seeds. The dashed vertical lines at $\sepratio = 1.1$ and $\sepratio = 1.3$ indicate the fallback and activation thresholds. All cell types with $\sepratio > 1.3$ show positive gain; all cell types with $\sepratio < 1.1$ show near-zero gain (|gain| $< 0.01$).}
\label{fig:sepratio}
\end{figure}

All five cell types with $\sepratio > 1.3$ show positive F1 gain (range $+0.023$ to $+0.777$), and all four cell types with $\sepratio < 1.1$ show $|\Delta\text{F1}| < 0.01$. The Spearman correlation between separability ratio and F1 gain across the nine cell types is $\rho = 0.81$ ($p = 0.009$). Crucially, the separability ratio is computed \emph{before} query inference, using only training-set embeddings---it requires no test-set evaluation or hyperparameter search.

The two highest-separability cases with near-zero gain (endothelial cell and pDC) share a common characteristic: scANVI baseline F1 is already high ($> 0.93$), leaving little room for improvement. This confirms that the separability ratio predicts the geometric capacity for rescue but not the amount of miscalibration in the scANVI classifier. A complete pre-deployment diagnostic combines both: rescue is expected when $\sepratio > 1.3$ \emph{and} scANVI baseline F1 (estimated on the validation set) is low ($< 0.75$).

\subsection{Fallback Guarantee on Low-Separability Cell Types}
\label{sec:fallback}

For the three cell types with $\sepratio < 1.1$ ($\varepsilon$-cells, $\gamma$-cells, $\beta$-cells), the prototype gate empties the rescue candidate set before marker verification is reached. As a consequence, Gate+Marker and scANVI produce identical predictions for all query cells in this regime, and the mean absolute F1 change across these cell types and all seeds is $0.003 \pm 0.007$. The small non-zero values reflect only numerical rounding. This guarantee is structural, not tuned: no test-set evaluation is needed to achieve it.

\subsection{Biological Plausibility of Rescued Cells}
\label{sec:markers}

The marker-gene signatures derived from labeled training cells are consistent with established biology. For ASDC, the top discriminative genes include \textit{AXL}, \textit{TCF4}, and \textit{LILRA4}, which are hallmark markers of the ASDC lineage \citep{villani2017single}. For cDC1, top markers include \textit{CLEC9A}, \textit{BATF3}, and \textit{ID2}, canonical cDC1 identity markers \citep{villani2017single}. For NCM, the signature is enriched for \textit{FCGR3A}, \textit{CX3CR1}, and \textit{RHOC}, known NCM markers. This alignment with published biology provides independent evidence that the rescued cells are genuine rare-class instances, not false positives.

\subsection{Framework Overview}
\label{sec:framework}

\cref{fig:framework} illustrates the complete scRareRefine workflow.

\begin{figure}[t]
\centering
\includegraphics[width=\linewidth]{scrarerefine_framework_prototype_gate_marker.pdf}
\caption{Overview of the scRareRefine pipeline. After scANVI training and query inference, the separability ratio $S$ is computed from labeled training cells. If $S \geq 1.1$, prototype scoring generates rescue candidates; the gate filters candidates by margin threshold; and marker-gene verification confirms biological plausibility. If $S < 1.1$, the pipeline falls back to scANVI predictions without modification. All reference quantities are derived from training data only; test-set labels are used solely for final evaluation.}
\label{fig:framework}
\end{figure}
